\begin{abstract}
\todo{Draft from the outline.}
\begin{itemize}
  \item \textbf{Computational problem:} solving the extended scalar-tensor-Gauss--Bonnet (ESTGB) field equations inside Bayesian inference, which needs $10^5$--$10^7$ evaluations, is a severe computational bottleneck~\citep{Reed2024}.
  \item \textbf{Physical scope:} a dual-branch surrogate framework for vacuum curvature-induced scalarized black holes, the 2D map $(\rh, \beta) \to (M, \Dch)$~\citep{Doneva2017_BH}, and matter-coupled scalarized neutron stars, the 3D map $(\beta, \lambda, \rhoc) \to (M, \Dch)$~\citep{Doneva2017_NS}.
  \item \textbf{Methodological novelty:} an exploratory data analysis of the $\nsEdaRawRows$ neutron-star configurations ($\nsEdaRows$ after cleaning) directly dictates feature extraction, target transformation, loss design and architecture selection ($H_1$, $H_2$, $H_3$).
  \item \textbf{Key results:} sub-percent mean absolute relative errors across the $\nsEdaDOrders$ orders of magnitude the dilaton charge $\Dch$ spans, speedups of $10^2$--$10^4\times$ (to be measured, Section~\ref{sec:speedup}), and unbiased MCMC recovery of the coupling parameters.
\end{itemize}
\end{abstract}
